\originalchapter{角动量与三维刚体}\label{ch:angular}
\originalsection{角动量及参考点}
关于固定惯性点 $O$, 质点角动量为 $\vct L_O=\vct r\times m\vct v$. 时间求导时 $\dot{\vct r}\times m\vct v=0$, 得 $\dot{\vct L}_O=\vct\tau_O$. 外力矩为零则角动量守恒. 单个中心力粒子即使线动量变化, 关于力中心的角动量仍守恒. 枢轴碰撞中支点可给强冲量但对支点无力矩, 这使角动量守恒优于线动量守恒.

对系统取质心分解, 得 $\vct L_O=\vct R\times M\vct V+\vct L_C$. 质心系中 $\sum m_i\vct r'_i=0$, 可以证明即使质心加速仍有 $\dot{\vct L}_C=\vct\tau_C$. 对任意移动参考点 $A$, 若角动量使用粒子的惯性系速度, 则 $\dot{\vct L}_A=\vct\tau_A-\vct v_A\times\vct P$. 省略修正项必须有固定点、质心或其他消项条件.

\originalsection{惯量张量与主轴}
三维转动中 $\vct L$ 一般不平行 $\vct\omega$. 将速度 $\vct v'=\vct\omega\times\vct r'$ 代入角动量, 用向量三重积得
\begin{equation}
L_i=\sum_jI_{ij}\omega_j,\qquad
I_{ij}=\int(r'^2\delta_{ij}-x'_ix'_j)\dd m.
\end{equation}
对角元如 $I_{xx}=\int(y'^2+z'^2)\dd m$, 非对角元 $I_{xy}=-\int x'y'\dd m$. 对称性使 $I$ 是实对称矩阵, 可选正交特征方向将它对角化. 这些方向称为主轴, 对应惯量 $I_1,I_2,I_3$. 主轴转动时角动量平行该轴, 一般转动则不平行. 动能为 $K_{\rm rot}=\vct\omega\cdot I\vct\omega/2$.

在随刚体转动的主轴系中, 任意向量的惯性导数等于该系分量导数加 $\vct\omega\times\vct A$. 应用于 $\vct L$ 得欧拉方程
\begin{align}
I_1\dot\omega_1+(I_3-I_2)\omega_2\omega_3&=\tau_1,\\
I_2\dot\omega_2+(I_1-I_3)\omega_3\omega_1&=\tau_2,\\
I_3\dot\omega_3+(I_2-I_1)\omega_1\omega_2&=\tau_3.
\end{align}
若 $I_1<I_2<I_3$, 无矩主轴转动的稳定性可由小扰动求出. 绕第 2 轴令 $\omega_2=\Omega$, 线性化第一、三式, 消去 $\omega_3$ 得
\begin{equation}
\ddot\omega_1=\frac{(I_2-I_3)(I_1-I_2)}{I_1I_3}\Omega^2\omega_1.
\end{equation}
右侧系数为正, 小扰动指数增长, 中间惯量轴不稳定. 绕最大或最小惯量轴则系数为负, 扰动振荡, 因而线性稳定. 惯量相等的退化情况须单独分析.

\originalsection{陀螺与冲量}
高速陀螺的自旋角动量大致沿轴, $\vct L\simeq I_s\omega_s\uv n$. 支点到质心为 $l\uv n$, 重力矩为 $l\uv n\times(-Mg\uv z)$. 若轴绕竖直方向以 $\vct\Omega=\Omega\uv z$ 稳定进动, $\dot{\vct L}=\vct\Omega\times\vct L$, 比较两式给 $\Omega\simeq Mgl/(I_s\omega_s)$. 进动方向取决于自旋符号. 该公式忽略进动自身角动量与章动, 限于快自旋、慢进动. 刚释放的陀螺一般先有章动, 不能说它从静止立刻严格水平进动.

短冲量 $\Delta\vct p$ 在 $\vct r$ 处作用, 角动量改变量为 $\Delta\vct L=\vct r\times\Delta\vct p$. 对自由质点这改变轨道角动量; 对非旋转杆或盘会从零建立转动; 对已有大自旋的盘, 小的垂直角冲量主要改变角动量方向, 随后可产生进动. 先算力矩与角动量, 再用主轴惯量求角速度, 能避免把角动量方向直接当成所有刚体的角速度方向.

欧拉盘近似中, 均匀薄盘半径 $a$, 盘面与水平夹角 $\alpha$, 质心位置近似不变. 无滑接触线是瞬时转轴, 总角速度沿接触点指向质心, 大小 $\Omega\sin\alpha$. 相应直径惯量为 $Ma^2/4$, 所以 $L=(Ma^2/4)\Omega\sin\alpha$. 角动量方向绕竖直转动的变化率大小为 $\Omega L\cos\alpha$, 重力与支持力的力偶为 $Mga\cos\alpha$, 因而 $\Omega^2=4g/(a\sin\alpha)$. $\alpha\to0$ 时该理想模型频率发散, 真实耗散与接触形变会使近似失效; 不能由此推出无限真实速度.
\originalsection{角冲量与受约束模型}
从固定惯性点的 $\dot{\vct L}_O=\vct\tau_O$ 积分, 得 $\Delta\vct L_O=\int_{t_1}^{t_2}\vct\tau_O\dd t$. 碰撞的“角动量守恒”要求所选点的外角冲量可忽略. 枢轴力可以令线动量改变但对枢轴矩为零; 重力矩在短碰撞中若有限, 其角冲量为较小的 $O(\Delta t)$ 项. 若冲击点在时间内明显移动, 不能简单把整个角冲量写成一个固定 $\vct r\times\Delta\vct p$.

均匀杆长 $l$, 一端无摩擦固定, 横向冲量 $J$ 作用在端点. 枢轴角冲量为 $lJ$, 惯量 $Ml^2/3$, 所以 $\omega=3J/(Ml)$, 质心速率 $V_C=\omega l/2=3J/(2M)$. 总外线冲量必须为 $MV_C$, 故支点冲量为 $J/2$、与敲击同向. 这具体表明“支点不做角冲量”不等于“不提供线冲量”. 同杆若自由, 质心速率为 $J/M$, 转动由质心角冲量 $Jl/2$ 与 $I_C=Ml^2/12$ 决定, $\omega=6J/(Ml)$. 改变支承就改变运动, 不能只换一个参考点而保持答案.

盘固定在可无摩擦转向的无质量轴上时, 沿轴作用的冲量可被轴承约束吸收, 垂直角冲量则改变盘的角动量. 若盘原本不转, $\Delta\vct L$ 是全部末角动量, 必须用与该方向对应的惯量求转速. 若盘原本沿对称轴快速自转, $L_s=I_s\omega_s$, 小的横向角冲量 $\Delta L_\perp$ 使方向偏角近似 $\Delta L_\perp/L_s$. 之后重力矩转动角动量方向形成进动, 同时可伴随章动. 从静止盘建立的翻转与自旋盘的轴向偏转来自同一角冲量定理, 但惯量、支承条件及后续力矩不同. 柱形陀螺的具体冲击例见期末第 5 题.

\originalsection{从倾斜轨道到力偶: 同一角动量原则}
教材第 22.3 节的模型链条强调方向而非只强调大小. 质量 $m$ 的小球用绳拴在中心, 在任意倾斜平面内作圆周运动, 无重力时绳张力为中心力, $\vct\tau_O=0$. 因而 $\vct L=mr^2\omega\uv n$ 是沿轨道平面法线的固定向量. 对球给横向冲量 $\vct J$, 则 $\vct L'=\vct L+\vct r\times\vct J$, 新轨道面垂直于 $\vct L'$. 这比“旋转变快了”更完整地描述敲击的效果.

将球换成无质量细杆两端的质量, 在位置 $\vct r_1,\vct r_2$ 给同时、等大反向的冲量 $\vct J,-\vct J$. 总线冲量为零, 角冲量为 $(\vct r_1-\vct r_2)\times\vct J$. 质心不加速, 但杆开始或改变转动. 更一般的刚轴上, 等大反向力 $\vct F,-\vct F$ 的力矩为 $(\vct r_1-\vct r_2)\times\vct F$: 它只取决于两作用点的相对位置, 不依赖选哪个参考点. 保持两力的间隔和方向而沿轴平移整个力偶, 对轴的净力矩不变; 单个力的力矩则没有这个性质.

轴上装盘但不自转时, 上述力偶建立横向 $\vct L$, 盘和轴直接绕相应横轴翻转. 盘已有很大的轴向自旋时, 同一个小横向角冲量主要把 $\vct L$ 偏转, 偏角 $\Delta L_\perp/L_s$. 力偶持续作用时 $\dot{\vct L}=\vct\tau$ 使大角动量向量的端点横向移动, 出现陀螺的转向; 要从角动量求轴的完整运动仍须惯量关系和支承条件. 因而“非旋转盘翻倒”与“旋转盘进动”可以沿同一模型链条理解, 不能把快陀螺近似套到静止盘上.

\originalsection{强制进动与支承载荷}
转盘上的陀螺由两个支座保持轴水平, 盘质心在两座中点, 各距 $d$. 盘自旋惯量 $I_s$, 速率 $\omega_s$, 平台绕竖直轴强制以 $\Omega$ 转动. 自旋角动量方向随平台改变, 所需水平力矩大小为 $I_s\omega_s\Omega$. 设两个向上支承力为 $F_A,F_B$, 选择 A 对应产生所需矩的一侧, 则
\[F_A+F_B=Mg,\qquad d(F_A-F_B)=I_s\omega_s\Omega.\]
所以 $F_A=(Mg+I_s\omega_s\Omega/d)/2$, $F_B=(Mg-I_s\omega_s\Omega/d)/2$. 若 $F_B=0$, 恰可撤去该支座形成自由进动; 若求得负值, 原来只向上托的支座不能维持该运动, 必须改为可拉住的轴承. 这把陀螺动力学变成可测的支承载荷, 不把全部重力平均分给两座.

磨坊轮是另一应用. 均匀轮 $M,b$, 质心沿半径 $R$ 水平圆轨道, 竖轴强制公转 $\Omega$, 水平无质量轮轴把轮心连到竖轴. 无滑要求自旋 $\omega_s=R\Omega/b$, 水平自旋角动量大小 $I_s\omega_s=MRb\Omega/2$. 取中央枢轴参考点, 公转角动量竖直且恒定, 水平自旋角动量随公转变化; 地面支持 $N$ 与重力的力矩差大小为 $R(N-Mg)$. 角动量方程给
\[R(N-Mg)=\frac12MRb\Omega^2,\qquad N=Mg+\frac12Mb\Omega^2.\]
向心加速度水平, 不会自行产生向下的竖直“向心力”. 额外压紧来自陀螺矩和轮轴的竖直约束力. 原教材要求轮以两倍重量压地, $N=2Mg$, 故 $\Omega=\sqrt{2g/b}$. 这是无滑、刚轴及均匀轮的理想模型.
